\section{Chapitre~\ref{sec:muscles}. La sortie vers les muscles}

La structure de la sortie --- unités motrices, marge de sécurité de la synapse, activation par taille, paire de muscles, boucle d'étirement et générateurs de rythme --- est mesurée directement; la condition de stabilité de la boucle à retard est un théorème; le modèle interne du corps est une hypothèse bien fondée; les chiffres des figures sont calculés avec le modèle du livre.

{\footnotesize
\setlength{\tabcolsep}{3pt}\renewcommand{\arraystretch}{1.15}
\begin{longtable}{>{\raggedright}p{0.5cm}>{\raggedright}p{4.0cm}>{\raggedright}p{5.6cm}>{\raggedright}p{2.3cm}>{\raggedright\arraybackslash}p{1.7cm}}
\toprule
N\textsuperscript{o} & Affirmation & Expérience et ce qui est mesuré & Source & Statut \\
\midrule\endfirsthead
\toprule
N\textsuperscript{o} & Affirmation & Expérience et ce qui est mesuré & Source & Statut \\
\midrule\endhead
1 & Unité motrice: un neurone \nt{ACET} commande un groupe de fibres, de quelques centaines à deux mille environ; dans les muscles à réglage fin, les unités sont plus petites & Muscles humains, comptage des axones moteurs et des fibres musculaires: en moyenne environ $340$ fibres par neurone dans un muscle interosseux de la main, environ $410$ dans le muscle brachio-radial, environ $1930$ dans le gastrocnémien médial. Le chapitre ne donne pas de chiffre précis pour les plus petites unités. & \cite{Feinstein1955mu} & mesuré \\
2 & La synapse sur le muscle libère plusieurs fois plus de neurotransmetteur qu'il n'en faut pour déclencher la fibre & Revue des mesures sur les synapses neuromusculaires: la marge de sécurité (rapport du neurotransmetteur libéré au seuil) vaut en général $3$--$5$ chez les mammifères adultes; chez l'homme, elle repose davantage sur les replis de la membrane du destinataire que sur l'importance de la libération. & \cite{WoodSlater2001} & mesuré \\
3 & La secousse~\eqref{eq:twitch} monte en un temps $T_c$ de l'ordre de $50\ms$ & Unités motrices isolées de la main humaine lors d'un effort volontaire, force moyennée sur les impulsions de l'unité: temps de montée de $30$--$100\ms$, inférieur à $70\ms$ pour $80\,\%$ des unités. La fonction $(t/T_c)e^{1-t/T_c}$ est une approximation tirée d'un modèle de pool d'unités. & \cite{MilnerBrown1973rec, Fuglevand1993} & mesuré \\
4 & Somme des secousses~\eqref{eq:forcesum}: à $5$, $15$ et $40$ imp./s, force de $0.2$--$1.1F_1$, $1.8$--$2.2F_1$ et environ $5.4F_1$ (fig.~\ref{fig:tetanus}) & Calcul avec \texttt{models/\allowbreak muscle\_\allowbreak models.py}, $T_c=50\ms$: $0.21$--$1.10$, $1.76$--$2.19$ et $5.28$--$5.49\,F_1$ sur les dernières $0.3\,\text{s}$. L'addition linéaire est une simplification: le modèle ne contient pas la saturation d'un vrai muscle. & \texttt{models/\allowbreak muscle\_\allowbreak models.py} & modèle \\
5 & Les unités s'activent par ordre de taille; les plus faibles sont cent fois plus faibles que les plus fortes & Racines ventrales du chat: sous une entrée réflexe croissante, les premiers neurones à décharger sont ceux dont les impulsions sont petites dans la racine, c'est-à-dire les petits. Muscle interosseux de la main humaine: force de secousse des unités de $0.1$ à $10$~g, croissant presque linéairement avec la force à laquelle l'unité s'active. & \cite{Henneman1957, MilnerBrown1973rec} & mesuré \\
6 & Le pas de force est proportionnel à la force, $\Delta F/F\approx q-1$~\eqref{eq:sizeprinciple}: virgule flottante & Déduit dans le chapitre de la progression géométrique des forces. Concorde avec l'expérience: aux efforts élevés, bien moins d'unités nouvelles s'activent pour un même incrément de force. & déduction du chapitre; \cite{MilnerBrown1973rec} & théorie \\
7 & La dispersion de la force croît proportionnellement à la force & Effort isométrique volontaire d'un muscle du pouce: l'écart type de la force croît linéairement avec la force moyenne; pour le même effort provoqué par stimulation électrique, la dispersion est constante. Modèle de pool: la croissance linéaire vient de l'activation des unités par ordre de force. & \cite{Jones2002sdn} & mesuré \\
8 & La fluidité des mouvements découle d'un bruit dépendant du signal & Modèle: la trajectoire qui minimise la dispersion du point d'arrivée sous un tel bruit reproduit les trajectoires mesurées de l'œil et du bras. Aucune expérience directe ne distingue cette cause des autres. & \cite{HarrisWolpert1998} & hypothèse \\
9 & La différence des forces d'une paire de muscles fixe le couple, leur somme la raideur~\eqref{eq:antagonists} & Théorie du contrôle de l'impédance par co-contraction. Bras humain: la raideur de chaque articulation, mesurée par de petites poussées, est à peu près proportionnelle au couple qui s'y exerce; différents rapports de co-contraction changent la forme et l'orientation de l'ellipse de raideur de la main. & \cite{Hogan1984, GomiOsu1998stiff} & mesuré \\
10 & Le capteur d'étirement excite son propre neurone \nt{ACET} et, par un intermédiaire inhibiteur, coupe l'antagoniste (fig.~\ref{fig:antagonists}) & Moelle épinière du chat: un seul intermédiaire, excité par les fibres des capteurs d'étirement, produit des potentiels inhibiteurs monosynaptiques dans les neurones \nt{ACET} du muscle antagoniste, délai synaptique $0.28$--$0.42\ms$. Le neurotransmetteur de l'intermédiaire (la glycine) n'a pas été déterminé dans cette expérience. & \cite{Jankowska1972Ia} & mesuré \\
11 & Retard de la boucle: $25$--$30\ms$ pour la boucle spinale, une fois et demie à trois fois plus par le cortex, $100$--$200\ms$ par l'œil & Bras humain, poussée sur l'articulation: réponse musculaire courte après $20$--$45\ms$, longue après $45$--$105\ms$, volontaire après $120$--$180\ms$. Décalage de la position visible du doigt pendant le mouvement: correction après $160\ms$ en moyenne. & \cite{Pruszynski2008rapid, SaundersKnill2003vis} & mesuré \\
12 & $\dd x/\dd t=-Gx(t-d)$ est stable pour $Gd<\pi/2$~\eqref{eq:delaystable}; à la limite, la fréquence est $1/(4d)$ & Théorème sur les racines d'une équation à retard. Vérification avec \texttt{models/\allowbreak muscle\_\allowbreak models.py}, $d=30\ms$: à $3\,\text{s}$, amplitude de $3\cdot10^{-6}$ pour $G=40$, $0.13$ pour $G=50$, $38$ pour $G=55$ par seconde (limite $52$). & \cite{Hayes1950} & théorie \\
13 & Tremblement physiologique de $8$--$12$~Hz; la boucle d'étirement est l'une de ses sources & Revue des mesures: un petit tremblement vers $10$~Hz existe chez les personnes en bonne santé; son origine est mixte: oscillations centrales, propriétés des décharges des unités, résonance mécanique et résonance de la boucle réflexe, et la part dominante varie selon les conditions. & \cite{McAuleyMarsden2000} & hypothèse \\
14 & Le cerveau prédit le résultat d'une commande grâce à un modèle interne du corps & Un sujet tient une charge entre deux doigts et bouge le bras: sous charge inertielle, visqueuse et élastique, la force de préhension varie en même temps que la charge, et non avec le retard de la boucle; la charge est donc prédite. Une revue rassemble ces expériences; on n'observe pas directement le modèle lui-même dans les neurones. & \cite{FlanaganWing1997grip, WolpertGhahramani2000} & hypothèse \\
15 & Le rythme de la marche, de la respiration, de la mastication est produit par des réseaux locaux sous entrée constante & Revue d'expériences: un système nerveux isolé (moelle épinière, ganglions) produit un motif moteur rythmique sans entrée rythmique et sans rétroaction des muscles. & \cite{MarderBucher2001} & mesuré \\
16 & L'inhibition mutuelle avec fatigue~\eqref{eq:halfcentre} donne un rythme de période $0.84\,\text{s}\approx2\tau_a$ (fig.~\ref{fig:halfcentre}) & Théorème: un réseau à inhibition mutuelle et adaptation sans état stable oscille nécessairement sous entrée constante. Calcul avec \texttt{models/\allowbreak muscle\_\allowbreak models.py} ($I=1$, $\beta=2$, $b=2$, $\tau=20\ms$, $\tau_a=0.4\,\text{s}$): période de $0.84\,\text{s}$. Que les vrais générateurs soient construits ainsi n'est pas démontré. & \cite{Matsuoka1985osc} & modèle \\
\bottomrule
\end{longtable}}
